\documentclass{article}

\usepackage{tabularx}
\usepackage{booktabs}

\title{Problem Statement and Goals\\\progname}

\author{\authname}

\date{}

\input{../Comments.text}
\input{../Common.text}

\begin{document}

\maketitle

\begin{table}[hp]
\caption{Revision History} \label{TblRevisionHistory}
\begin{tabularx}{\textwidth}{l>{\raggedright\arraybackslash}p{0.3\textwidth}>{\raggedright\arraybackslash}X}
\toprule
\textbf{Date} & \textbf{Developer(s)} & \textbf{Change}\\
\midrule
2026-09-25 & Anam Khan, Haniya Kashif, Vivian Nguyen, Kelsey Rough, Hareem Salman
& Initial version of the problem statement and goals: problem, inputs and
outputs, stakeholders, environment, goals, stretch goals, custom documents, and
reflection\\
\bottomrule
\end{tabularx}
\end{table}

\section{Problem Statement}

% \wss{You should check your problem statement with the
% \href{https://smiths.github.io/capTemplate/Checklists/ProbState-Checklist.pdf}
% {problem statement checklist}}. 

% \wss{You can change the section headings, as long as you include the required
% information.}

In the laboratory, Doppler ultrasound measurements can be used to study blood flow in arteries such as the brachial and femoral arteries. Previously, the laboratory used a Multigon system to process the audio output from an ultrasound machine and convert the Doppler signal into a voltage signal that could be continuously recorded and analyzed using LabChart.

The existing Multigon system is no longer functional, creating a need for a modern replacement that will take an input of the ultrasound audio waves and create an output voltage signal in real time. Without a functioning system, researchers cannot readily convert the Doppler ultrasound output into a continuous measure of MBV (Mean Blood Velocity) using the laboratory's existing recording workflow, creating an obstacle in their path to collect data and continue their research.

\subsection{Problem}

% \wss{What problem are you trying to solve?  Why is it important? Do not talk
% about how you will solve the problem, only what the problem is.  AI is rarely
% part of the problem.  The problem is ``what'' you want to do, not ``how'' you
% want to do it.  The problem statement should be written in a way that is
% understandable to a non-technical audience.}

We want to provide a system that will receive the Doppler audio output from an ultrasound machine, process the signal, and generate a continuous voltage output that can be recorded by LabChart for subsequent analysis of mean blood velocity.

\subsection{Inputs and Outputs}

The primary input to the system is the Doppler ultrasound audio signal representing blood flow within an artery.
 
The primary output is a continuous signal representing mean blood velocity that can be recorded and analyzed using LabChart. The system should also provide sufficient information to allow researchers to verify that the resulting measurements are reasonable and reliable.

\subsection{Stakeholders}

The primary stakeholders are the researchers and laboratory staff who require MBV measurements for physiological research. The project supervisors and biomedical researchers are also stakeholders, as they will provide domain knowledge, define research requirements, and evaluate the system's results. Ultrasound operators are stakeholders because the system must work with the output produced by the laboratory's ultrasound equipment. Future research participants are indirect stakeholders because the system may eventually be used to obtain measurements during human research studies.


\subsection{Environment}

% \wss{Hardware and Software Environment for the users.  The developer environment
% is summarized as part of the developer plan.}

The system will operate in a laboratory environment alongside an existing Doppler ultrasound system and LabChart data acquisition software. Researchers will use the system to obtain continuous MBV measurements during physiological experiments. The system will require appropriate hardware for receiving the ultrasound audio output and producing a signal that can be recorded by LabChart.

\section{Goals}

The primary goal of this project is to develop and validate a digital signal processing system that converts Doppler ultrasound audio into a continuous estimate of mean blood velocity (MBV), with an output suitable for recording and analysis in LabChart. The system will provide researchers with a modern and reliable alternative to the previously used Multigon system.
 
The system goals are to:
 
\begin{enumerate}
    \item \textbf{Define hardware and signal-conversion requirements:} The system shall identify and document the hardware interface requirements between the all the equipment, including input/output signal characteristics, sampling requirements, voltage ranges, impedance, and required signal conversion. These specifications will establish the hardware and software requirements needed to convert the ultrasound signal into a LabChart-compatible representation of mean blood velocity.
    \item \textbf{Acquire and process Doppler ultrasound signals:} The system will acquire the audio output from the ultrasound equipment and process the signal continuously with a sampling rate and frequency range appropriate for the input signal.
    \item \textbf{Extract blood-flow information from the Doppler signal:} The system shall apply and evaluate signal-processing techniques, including frequency-domain analysis, to identify Doppler frequency components associated with blood flow.
    \item \textbf{Estimate mean blood velocity:} The system should convert the extracted Doppler frequency information into a continuous estimate of mean blood velocity using an established Doppler velocity estimation method.
    \item \textbf{Provide a LabChart-compatible output:} The system shall generate a continuous output signal representing mean blood velocity that satisfies the electrical and signal requirements of the labs LabChart setup.
    \item \textbf{Validate measurement accuracy:} Will evaluate the system against an appropriate reference measurement or established method to quantify the accuracy and reliability of its mean blood velocity estimates.
    \item \textbf{Provide a usable and reproducible research tool:} The system will include documentation describing its operation, signal-processing pipeline, configuration, validation results, and known limitations so that researchers can reproduce and use the system in future studies.
\end{enumerate}

\section{Stretch Goals}

% \wss{Goals you hope to get to, but not necessary.  They are also helpful if the
% scope of the original project needs to be expanded.}

Additional goals we'd hope to implement in CADENCE include but are not limited to the following:
\begin{enumerate}
    \item \textbf{Extend physiological measurements:} Using the same input signal from the ultrasound, we can apply different signal transformation algorithms to derive metrics other than just MBV, such as Peak systolic velocity (PSV), End-diastolic velocity (EDV), Pulsatility Index (PI), and Resistance Index (RI). This would increase the usefulness of CADENCE beyond mean blood velocity alone and provide researchers with additional measures for physiological analysis.
    \item \textbf{Real-time signal visualization:} Alongside the LabChart output, provide a real-time spectrogram (a color-coded time-frequency display) to provide a more insightful diagnostic trail for research purposes, allowing them to visualize the frequency content and temporal characteristics of the Doppler signal during data collection. This would provide an additional means of assessing signal quality and interpreting the system's processing behaviour. It can also then be used by users who do not necessarily have LabChart.
    \item \textbf{Modular signal-processing framework:} Design the signal-processing pipeline so that additional signal-processing algorithms and physiological measurements can be incorporated without redesigning the entire system. This would allow CADENCE to be extended for future research applications.
    \item \textbf{AI-assisted signal analysis:} Explore the use of artificial intelligence or machine learning to analyze the processed Doppler signal and identify notable patterns or events in blood-flow measurements. For example, the system could detect significant changes or spikes in blood velocity and provide contextual information about the observed signal pattern. This could provide researchers with an additional layer of automated analysis and help identify features of interest for further investigation.
\end{enumerate}

\section{Custom Documents (CDs)}

% \wss{For CAS 741: State whether the project is a research project. This
% designation, with the approval (or request) of the instructor, can be modified
% over the course of the term.}

% \wss{For SE Capstone: List your custom documents (CDs).  Potential CDs include
% usability testing, code walkthroughs, user documentation, formal proof,
% GenderMag personas, Design Thinking, etc.  (The full list is on the course
% outline and in Lecture 02.) Normally the number of CDs will be two (one per
% term).  Approval of the CDs will be part of the discussion with the instructor
% for approving the project. The CDs, with the approval (or request) of the
% instructor, can be modified over the course of the term.}

For the Fall term, we have chosen to complete a User/Stakeholder/Client
Interview Report. Our project is being developed for the Vascular Dynamics Lab,
and the researchers and graduate students there are the people who will
actually be using our system. Many of the lab's requirements come from how they
currently use the Multigon, such as its velocity ranges, offset, calibration
signal, and output voltage, and a lot of this knowledge isn't formally
documented anywhere. Because nobody on our team has a background in the medical
field, interviewing the lab and our supervisors will help us capture these
requirements accurately and bridge the knowledge gap between our team and the
lab. This will also help us make sure that our system fits into the lab's
current requirements and how they hope to use the system.


For the Winter term, we have chosen to complete a Performance Report. Since our
system needs to process the Doppler signal and produce a continuous output in
real-time, it is essential that the system's output be both precise and
accurate. The delay between the incoming Doppler signal and the output voltage
recorded in LabChart directly affects whether the lab can use our measurements
for research. A performance report will allow us to formally measure and
analyze the latency and accuracy of our final system on its hardware, and
determine whether or not it meets the project requirements and the lab's needs.

\newpage{}

\section*{Appendix --- Reflection}

% \wss{Not required for CAS 741}

% \input{../Reflection.text}

\begin{enumerate}
    \item What went well while writing this deliverable? 

    For this deliverable, we found that identifying the core problem and user
    needs was straightforward. The Vascular Dynamics Lab had an existing
    Multigon machine that is broken and discontinued, and they need a
    replacement. They told us clearly that they were looking for a system that
    takes the ultrasound's Doppler audio and outputs a continuous mean blood
    velocity voltage signal for their LabChart software to interpret. The lab
    has been incredibly helpful in providing us with manuals, test data, and
    time in the lab to better understand their system and how it works. They
    also gave us ideas for stretch goals and requirements, and helped us
    understand their primary user needs. We were also able to get help from
    Dr. Von Mohrenschildt, who advised us on what information we need to gather
    about the system's electronics and existing signals, which helped us
    understand what approach we need to take with the project.

    Within our team, we split up the work evenly, each taking on parts of the
    problem statement, goals, and project research. When it came to determining
    all of the technical details we needed, we met to research together, each
    taking on specific elements or interfaces within the system. We also shared
    responsibilities such as contacting supervisors and the lab to arrange
    meetings, tracking the administrative tasks we needed to do, and organizing
    our shared resources as a group.

    We found that meeting with the lab and our supervisor early was one of our
    best decisions. Seeing the equipment in person changed how we wrote the
    problem statement because instead of describing a generic signal processing
    tool, we were able to focus it on the specific input and output the lab
    actually needs. It also helped us identify the lab's needs, along with the
    information and context we needed early on, which gave us what we needed to
    accurately plan our goals and development.

    \item What pain points did you experience during this deliverable, and how
    did you resolve them?

    The biggest pain point we found for this deliverable was the number of
    unknowns about the existing system our project aims to replace. Because it's
    legacy hardware, there's little documentation for operating it, much less
    technical details such as current and voltage draw, interface impedances,
    signal formats and phase shifts, data bit depth, and other key details we
    need to accurately determine our next steps. We addressed a lot of this by
    going into the lab a couple of times to see the machine, watch some tests
    being done, and see how the output looks. We haven't resolved all the
    unknowns yet, so we wrote our goals around what the system needs to deliver
    rather than around specs we haven't confirmed, and we will be updating these
    as soon as we have confirmed answers. We're returning to the lab next week
    to measure the remaining unknowns with a multimeter and oscilloscope, and
    we'll update the documents based on what we find. A lot of this project is
    new to our team, so we've had to do a lot of research and learning as we
    go, and looking back we recognize that we may have also benefitted from
    putting together a list of specific measurements to take before our first
    lab visit instead of after.

    One of the other pain points we experienced was the difference in domain
    knowledge on both sides. The Vascular Dynamics Lab is an amazing resource on
    the physiology side, but they don't have the technical engineering
    background that we do. On the other hand, we've had a lot to learn about the
    medical side of this project and its implications for our design. To bridge
    this gap, when we visited the lab, we asked them to demonstrate to us how
    they use the system, and we took notes on how it works, what readings they
    get, and what may be required for us to provide. We also had a gap between
    our current technical knowledge and what the project requires. Dr. Von
    Mohrenschildt was a great help with this. When we met with him, he explained
    key concepts and helped us see what we need for the project and how to
    approach it. The main thing we learned is that asking specific questions got
    us much better answers than general ones, so we plan to come to future
    meetings with questions prepared ahead of time.

    \item How did you and your team adjust the scope of your goals to ensure
    they are suitable for a Capstone project (not overly ambitious but also of
    appropriate complexity for a senior design project)?

    As our team looked further into the system we need to develop, it became
    clear that we needed to adjust our goals to make sure they are feasible and
    suitable for this capstone. Initially, we had ideas of providing other
    outputs, such as peak systolic velocity, end-diastolic velocity, pulsatility
    index, and resistance index, along with a real-time spectrogram display.
    However, from being in the lab, examining the equipment, and speaking with
    the lab group, it became clear that the core problem is complex enough on
    its own. Most of these extra metrics can be derived from the same velocity
    data, but each one would need its own implementation, testing, and
    verification. That would pull time and effort away from getting the core
    mean blood velocity output right. The lab also has ways of interpreting
    these already, so the addition of these would add little to their current
    research. Because this system will be used for medical research, an
    accurate and reliable core output matters far more than extra features. We
    adjusted our goals to prioritize system accuracy and reliability, along with
    clear documentation and verification of correctness, so the system is fit
    for research use and easy for the lab to operate. We haven't decided against
    the additional metrics entirely, and we do have some stretch goals that we
    aim to complete if the core output is validated with time to spare. We
    believe this scope is appropriately challenging, because it involves
    real-time signal processing, hardware interfacing, and validation against
    real laboratory equipment, which are all relevant to the workforce and also
    allow us to apply concepts learned in SFWRENG 3MX3, 3DX4, and 4AA4.
\end{enumerate}  

\end{document}